\section{Introduction}

Exact compilation is usually posed as optimization over a sufficiently expressive search space. A complementary question is whether that space admits a smaller exact normal form. Such a result can separate representational necessity from implementation performance and can expose the precise obstruction to a stronger restriction.

We study this question for a fixed Tag-and-Restore construction derived from a stabilizer-formalism block-encoding method \cite{schillo2026blockencoding}. An instance contains six nonidentity Pauli targets arranged into three ordered blocks. Each block uses an anticommuting auxiliary-frame pair, and the three blocks share one label operator. The declared objective counts auxiliary support and factored correction support with fixed coefficients. It is a logical structural coordinate, not a circuit-depth, fault-tolerant, or hardware cost model.

\paragraph{Main results and assumptions.}
All statements use exactly this six-target, three-block grammar and objective. First, we prove that every optimum has a representative in which each frame Pauli has support at most two. Second, an exact two-qubit counterexample proves that support one is not uniformly sufficient, so the uniform frame-support threshold is exactly two. Third, we count the ordered anticommuting support-two pairs and obtain a direct exact $O(n^9)$ word-RAM algorithm. Fourth, finite checks corroborate the local lemmas and exact implementation. Finally, a pre-specified runtime study fails to show implementation benefit: the direct algorithm times out on all six full-subject comparisons.

The distinction between theorem and performance evidence is central. The normal form is an exact mathematical statement. The asymptotic result is an upper bound for a new direct algorithm, not a lower bound or a speed claim. The runtime study is retained as adverse evidence rather than converted into a positive result. No statement here implies physical-resource advantage or optimality for another grammar, objective, or compiler.
